\documentclass[12pt, a4paper]{article}

% ---- Ajustes de Idioma e Fontes Modernas -------------------
\usepackage[utf8]{inputenc}
\usepackage[T1]{fontenc}
\usepackage[brazil]{babel}
\usepackage{mathpazo}
\usepackage[scaled=0.95]{helvet}
\usepackage{courier}

% ---- Pacotes de Layout e Design ----------------------------
\usepackage{geometry}
\geometry{a4paper, top=3cm, bottom=2cm, left=3cm, right=2cm}
\usepackage{setspace}
\onehalfspacing
\usepackage{indentfirst}
\usepackage{microtype}

% ---- Cores e Hiperlinks ------------------------------------
\usepackage{xcolor}
\definecolor{brandblue}{RGB}{44, 62, 80}
\definecolor{softgray}{RGB}{245, 245, 245}
\usepackage[colorlinks=true, allcolors=brandblue]{hyperref}

% ---- Elementos Matemáticos e Gráficos ----------------------
\usepackage{amsmath, amssymb, amsfonts}
\usepackage{graphicx}
\usepackage{booktabs}
\usepackage{caption}
\usepackage{subcaption}
\usepackage{float}

% ---- Customização de Títulos -------------------------------
\usepackage{titlesec}
\titleformat{\section}{\color{brandblue}\normalfont\Large\bfseries}{\thesection}{1em}{}
\titleformat{\subsection}{\color{brandblue}\normalfont\large\bfseries}{\thesubsection}{1em}{}

% =============================================================
\begin{document}
% =============================================================

% ---- CAPA ---------------------------------------------------
\begin{titlepage}
    \begin{flushleft}
        \Large \textbf{Instituto Federal do Ceará} \\
        \large Mestrado em Ciência da Computação \\
        \small Disciplina: Reconhecimento de Padrões
    \end{flushleft}
    \vspace{4cm}
    \begin{center}
        \hrule\vspace{0.5cm}
        {\huge \textbf{\color{brandblue} Classificadores Bayesianos}}\\
        \vspace{0.2cm}
        {\Large Trabalho 02: Gaussiano Multivariado}
        \vspace{0.5cm}\hrule
    \end{center}
    \vspace{3cm}
    \begin{flushright}
        \large \textbf{Raquel Maciel Coelho de Sousa}\\
        \small Professor: Dr. Ajalmar Rocha Neto
    \end{flushright}
    \vfill
    \begin{center}\small Fortaleza, 11 de maio de 2026 \end{center}
\end{titlepage}

\newpage
{\hypersetup{linkcolor=black}\tableofcontents}
\newpage

% =============================================================
\section{Introdução}
% =============================================================

Este relatório descreve a implementação e avaliação do
\textbf{Classificador Bayesiano Gaussiano Multivariado}, desenvolvido
como Trabalho~02 da disciplina de Reconhecimento de Padrões. O modelo
foi treinado e testado sobre cinco conjuntos de dados: Flor Íris,
Coluna Vertebral, Breast Cancer Wisconsin, Dermatology e Artificial~I, e seus resultados foram comparados
com os classificadores KNN e DMC.

% =============================================================
\section{Fundamentação Teórica}
% =============================================================

\subsection{Teorema de Bayes}

O classificador fundamenta-se no Teorema de Bayes:

\begin{equation}
    P(C_i \mid \mathbf{x}) =
    \frac{p(\mathbf{x} \mid C_i)\,P(C_i)}{p(\mathbf{x})}
\end{equation}

\noindent onde $P(C_i)$ é a probabilidade \textit{a priori} da classe,
$p(\mathbf{x} \mid C_i)$ é a verossimilhança e
$p(\mathbf{x}) = \sum_k p(\mathbf{x} \mid C_k)\,P(C_k)$ é a evidência.
A decisão final é $\hat{C} = \arg\max_{C_i} P(C_i \mid \mathbf{x})$.

\subsection{Distribuição Gaussiana Multivariada}

A verossimilhança é modelada pela gaussiana multivariada de dimensão $d$:

\begin{equation}
    p(\mathbf{x}) =
    \frac{1}{(2\pi)^{d/2}|\Sigma|^{1/2}}
    \exp\!\left\{
        -\tfrac{1}{2}(\mathbf{x}-\boldsymbol{\mu})^T
        \Sigma^{-1}(\mathbf{x}-\boldsymbol{\mu})
    \right\}
    \label{eq:gaussiana}
\end{equation}

\noindent onde $\boldsymbol{\mu}_i$ e $\Sigma_i$ são estimados
\textbf{independentemente para cada classe} a partir dos dados de treino,
sem simplificações.

\subsection{Matriz de Covariância}

\begin{equation}
    \Sigma =
    \begin{bmatrix}
        \sigma_{11}^2 & \sigma_{12} & \cdots & \sigma_{1d}\\
        \sigma_{21} & \sigma_{22}^2 & \cdots & \sigma_{2d}\\
        \vdots & \vdots & \ddots & \vdots\\
        \sigma_{d1} & \sigma_{d2} & \cdots & \sigma_{dd}^2
    \end{bmatrix}
\end{equation}

A diagonal contém as variâncias individuais; fora da diagonal estão
as covariâncias, que capturam correlações entre pares de atributos.
A simetria ($\Sigma_{ij}=\Sigma_{ji}$) é explorada para reduzir
cálculos: apenas o triângulo superior é computado.

\subsection{Regularização de Tikhonov}

Quando um atributo tem variância zero dentro de uma classe (comum em
datasets com atributos binários e classes pequenas), $\Sigma$ torna-se
singular e não possui inversa. A solução adotada é somar
$\varepsilon\mathbf{I}$ à diagonal antes de inverter:

\begin{equation}
    \Sigma_{\text{reg}} = \Sigma + \varepsilon\mathbf{I},\quad
    \varepsilon = 10^{-6}
\end{equation}

\subsection{Classificadores de Referência}

\textbf{DMC} — no treino calcula o centroide de cada classe
$\boldsymbol{c}_i = \frac{1}{|C_i|}\sum_{\mathbf{x}\in C_i}\mathbf{x}$;
na classificação escolhe a classe de menor distância euclidiana.

\textbf{KNN} ($k=5$) — não possui treino explícito; classifica pelo
voto majoritário entre os $k$ vizinhos mais próximos no conjunto de
treino.

\subsection{Escolha do Par de Atributos para Visualização}

O par foi selecionado maximizando o critério de Fisher :

\begin{equation}
    \text{score}(a_1,a_2) =
    \text{Fisher}(a_1)+\text{Fisher}(a_2)
\end{equation}

\noindent onde $\text{Fisher}(a) = \frac{\text{Var}(\mu_1,\ldots,\mu_K)}
{\text{média}(\sigma^2_1,\ldots,\sigma^2_K)}$ mede o quanto as classes
se afastam em relação ao seu espalhamento interno.

% =============================================================
\section{Conjuntos de Dados}
% =============================================================

\subsection{Flor Íris}
150 amostras, 4 atributos contínuos, 3 classes balanceadas.
Dataset limpo, sem valores faltantes.

\subsection{Coluna Vertebral}
310 amostras, 6 atributos biomecânicos contínuos, 3 classes
ortopédicas. Completamente numérico.

\subsection{Breast Cancer Wisconsin}
569 amostras, 30 atributos numéricos extraídos de imagens digitais,
2 classes. Foi utilizada a versão \textit{Diagnostic} (UCI id=17),
completamente numérica. A versão original (id=14) foi descartada
por conter atributos categóricos que exigiriam imputação complexa.

\subsection{Dermatology}
366 amostras, 34 atributos, 6 classes. Valores faltantes foram
imputados pela média da coluna.

\subsection{Artificial I}
Dataset sintético gerado com 3 classes gaussianas 2D
($n=50$ por classe, semente $42$):
\begin{itemize}
    \item \textbf{Círculo:} $\boldsymbol{\mu}=(1.0,\,4.0)$, $\sigma=0.5$
    \item \textbf{Triângulo:} $\boldsymbol{\mu}=(4.0,\,4.0)$, $\sigma=0.5$
    \item \textbf{Estrela:} $\boldsymbol{\mu}=(2.0,\,2.0)$, $\sigma=0.5$
\end{itemize}

% =============================================================
\section{Metodologia}
% =============================================================

Foram realizadas \textbf{25 realizações} independentes por dataset.
Em cada realização os dados foram embaralhados e divididos em 80\%
treino e 20\% teste. Os quatro classificadores foram avaliados sobre
a mesma divisão. Ao final:

\begin{equation}
    \bar{A} = \frac{1}{N}\sum_{n=1}^N A_n
    \qquad
    \sigma_A = \sqrt{\frac{1}{N}\sum_{n=1}^N(A_n-\bar{A})^2}
\end{equation}

A matriz de confusão foi extraída da realização com acurácia
\textbf{mais próxima da média}, por ser a mais representativa.

% =============================================================
\section{Resultados e Discussão}
% =============================================================

\subsection{Acurácia Média e Desvio Padrão}

\begin{table}[H]
    \centering
    \caption{Resultados experimentais --- 25 realizações, divisão 80/20.}
    \label{tab:results}
    \begin{tabular}{@{}lcccc@{}}
        \toprule
        \textbf{Base de Dados} & \textbf{Bayesiano} & \textbf{KNN ($k=5$)} & \textbf{DMC} \\
        \midrule
        Íris             & $97.5 \pm 3.0$ & $96.3 \pm 2.9$ & $90.9 \pm 4.9$ \\
        Artificial I     & $99.2 \pm 1.4$ & $99.2 \pm 1.7$ & $99.2 \pm 1.4$ \\
        Coluna Vertebral & $83.9 \pm 4.5$ & $82.3 \pm 3.3$ & $75.7 \pm 5.5$ \\
        Breast Cancer    & $88.4 \pm 3.0$ & $93.6 \pm 2.2$ & $87.9 \pm 2.7$ \\
        Dermatology      & $88.6 \pm 3.8$ & $86.2 \pm 2.9$ & $50.4 \pm 6.2$ \\
        \bottomrule
    \end{tabular}
\end{table}

\subsection{Matrizes de Confusão}
Matriz de confusão de cada dataset com o classificador Bayesiano Multivariado. Buscando a melhor representatividade, foram utilizadas os registros da realização com a taxa de acerto mais próxima da acurácia para montar cada matriz.

\begin{figure}[H]
    \centering
    \includegraphics[width=0.7\textwidth]{iris-matrix-confusao.png}
    \caption{Matriz de confusão --- Íris.}
    \label{fig:conf_multi}
\end{figure}

\begin{figure}[H]
    \centering
    \includegraphics[width=0.7\textwidth]{vertebral-matrix-confusao.png}
    \caption{Matriz de confusão --- Vertebral Column.}
    \label{fig:conf_multi}
\end{figure}

\begin{figure}[H]
    \centering
    \includegraphics[width=0.7\textwidth]{breast-matrix-confusao.png}
    \caption{Matriz de confusão --- Breast Cancer.}
    \label{fig:conf_multi}
\end{figure}

\begin{figure}[H]
    \centering
    \includegraphics[width=0.7\textwidth]{dermatology-matrix-confusao.png}
    \caption{Matriz de confusão --- Dermatology.}
    \label{fig:conf_multi}
\end{figure}

\begin{figure}[H]
    \centering
    \includegraphics[width=0.7\textwidth]{artificial-matrix-confusao.png}
    \caption{Matrix de confusão --- Artificial I.}
    \label{fig:conf_multi}
\end{figure}

\subsection{Superfícies de Decisão}
Superfícies de decisão do Bayesiano Multivariado com gaussianas sobrepostas. Treino~($\bullet$), teste ($\times$).

\begin{figure}[H]
\centering
\includegraphics[width=1\textwidth]{iris-superficie-decisao.png}
\caption{Superfície de decisão --- Íris}
\end{figure}

\begin{figure}[H]
\centering
\includegraphics[width=1\textwidth]{artificial-superficie-decisao.png}
\caption{Superfície de decisão --- Artificial I}
\end{figure}

\begin{figure}[H]
\centering
\includegraphics[width=1\textwidth]{vertebral-superficie-decisao.png}
\caption{Superfície de decisão --- Vertebral Column.}
\end{figure}

\subsection{Discussão}

O \textbf{Artificial~I} atingiu acurácia próxima de 100\% em todos os
classificadores, confirmando que classes gaussianas bem separadas
representam o cenário ideal para modelos probabilísticos.

Na \textbf{Íris}, o Bayesiano Multivariado ($97.5\%$) supera o DMC
($90.9\%$) e iguala o KNN ($96.3\%$), beneficiando-se da modelagem
das correlações entre atributos via $\Sigma_i$.

No \textbf{Breast Cancer}, o KNN ($93.6\%$)
supera o Bayesiano Multivariado ($88.4\%$). Isso sugere que, com 30
atributos, a estimativa da matriz de covariância $30\times30$ a partir
de apenas $\sim$455 amostras de treino introduz ruído prejudicial ---
um fenômeno conhecido como \textit{curse of dimensionality}.

No \textbf{Dermatology}, o DMC apresenta queda expressiva ($50.4\%$),
evidenciando que centroides são insuficientes para separar 6 classes
com distribuições complexas. O Bayesiano Multivariado ($88.6\%$)
destaca-se como o melhor classificador neste dataset.

% =============================================================
\section{Conclusão}
% =============================================================

O Bayesiano Gaussiano Multivariado mostrou-se competitivo em todos os
datasets, destacando-se especialmente no Dermatology. Sua principal
limitação aparece em alta dimensionalidade (Breast Cancer), onde a
estimativa de $\Sigma$ torna-se instável. Os desafios práticos
encontrados foram:

\begin{itemize}
    \item \textbf{Overflow no expoente:} resolvido com \textit{clipping}
          ao intervalo $[-500,500]$;
    \item \textbf{Matriz singular:} resolvido com regularização de
          Tikhonov ($\varepsilon=10^{-6}$);
    \item \textbf{Divisão treino/teste:} garantido mínimo de 1 amostra
          de teste por realização.
\end{itemize}

\newpage
\begin{thebibliography}{9}
    \bibitem{duda} DUDA, R. O.; HART, P. E.; STORK, D. G\@.
        \textbf{Pattern Classification}. 2.~ed. Wiley-Interscience, 2001.
    \bibitem{bishop} BISHOP, C. M\@.
        \textbf{Pattern Recognition and Machine Learning}. Springer, 2006.
    \bibitem{uci} DUA, D.; GRAFF, C\@. UCI Machine Learning Repository.
        Disponível em: \url{https://archive.ics.uci.edu/ml}. Acesso em: \today.
\end{thebibliography}

\end{document}